都满足上述条件。在这一语境里存在\trMoser{}引理的一个弱类比，可陈述如下。设 $M$、$N$
为\trCorneredBoundary{}的\trExactSymplectic\trManifold。一个\trEmbedding{}$\rho\colon M\to N$ 称为%
\textbf{\trExact\trConformallySymplectic{}\trEmbedding}（exact conformally symplectic embedding），
若对某个 $c\in\RR$ 有 $\rho^*\theta_N=e^{-c}\theta_M+d(\text{某个函数})$。那么，
\begin{lemma}\label{thm:moser-deformation-embedding}
若两个\trCorneredBoundary{}的\trExactSymplectic\trManifold{}可经\trDeformation{}互相变入对方，
则存在一个\trExact\trConformallySymplectic{}\trEmbedding，把其中一个\trEmbedding{}到另一个之中。
\end{lemma}
\begin{proof}
取一个\trDeformation，并如标准的\trMoser{}论证那样，考察由 $\partial\theta^s/\partial s$ 得到的%
依赖于 $s$ 的\trVectorField。它未必沿\trBoundaryFace\trIncoming，但这一缺陷可以通过添上 $X_{\theta^s}$
的一个大倍数来补救。把所得的\trVectorField{}积分，便给出所要的\trEmbedding。
\end{proof}
还能说得稍微精确些。考虑我们刚刚构造的那个\trEmbedding{}$\rho$ 在全体%
\trExact\trConformallySymplectic{}\trEmbedding{}所成空间中的\trIsotopy{}类。它不依赖于任何选择；在%
\trDeformation{}的保持端点的\trHomotopy{}下不变；并且与\trDeformation{}的\trComposition{}\trCompatible。
特别地，由此可知：若 $\widetilde{\rho}$ 是把该\trDeformation{}倒着跑一遍所得的\trEmbedding，则%
$\widetilde{\rho}\circ\rho$ \trIsotopy{}于恒等。
\paragraph{\trContactTypeBoundary{}的情形。}\label{sub:contact-type-boundary}
作为对上文讨论的补充，我们现在来看一类更为人熟悉的带边\trSymplectic\trManifold。一些奠基性文献是
[42]、[100]。一个\textbf{\trContactTypeBoundary{}的\trExactSymplectic\trManifold}（exact symplectic
manifold with contact type boundary）是一个\trCompact\trManifoldWithBoundary（但无\trCornerPoint）%
$M$，配有\trSymplecticForm{}$\varpi_M$ 与 1-\trForm{}$\theta_M$，满足 $d\theta_M=\varpi_M$，使得%
$X_{\theta_M}$ 沿 $\partial M$ \trIncoming。
一个基本构造以\trCanonical{}的方式粘上一个无穷锥，把这样的 $M$ 放大成非紧的%
\trExactSymplectic\trManifold：
\begin{equation}\label{eq:cone-completion}
(\widehat{M},\varpi_{\widehat{M}},\theta_{\widehat{M}})
=(M,\varpi_M,\theta_M)\cup_{\partial M}
\bigl(\RR^+\times\partial M,\;d\bigl(e^r(\theta_M|\partial M)\bigr),\;e^r(\theta_M|\partial M)\bigr),
\end{equation}
其中 $r$ 是锥的\trRadial{}变量（本书全书的约定是 $\RR^\pm\subset\RR$ 取闭，即包含 $0$）。沿“接缝”的%
可微结构由两侧的\trLiouville\trFlow{}所提供的 $\partial M$ 的\trCollar{}邻域来定义。于是，$M$ 上的%
负\trLiouville{}场 $X_{\theta_M}$ 与 $\RR^+\times\partial M$ 上的 $-\partial_r$ \trGluing{}在一起，
给出 $\widehat{M}$ 上的负\trLiouville{}场。类似地注意，存在唯一的（在\trGerm{}的意义下）函数%
$h_M\colon U\to(0;1]$，定义于 $\partial M$ 的某个邻域 $U\subset M$ 之上，使得
\begin{equation}\label{eq:contact-gluing-function}
h_M^{-1}(1)=\partial M \qquad\text{且}\qquad X_{\theta_M}.h_M=-h_M.
\end{equation}
这可以\trSmooth{}地延拓为 $U\cup(\RR^+\times\partial M)\subset\widehat{M}$ 上的函数%
$h_{\widehat{M}}$，办法是在锥上令 $h_{\widehat{M}}(r,y)=e^r$。
在\eqref{eq:cone-completion} 型的两个\trManifold{}之间，一个恰当的\trIsomorphism{}概念由那些满足%
$\widehat{\varphi}^*\theta_{\widehat{N}}-\theta_{\widehat{M}}=d(\text{某个函数})$ 的%
\trDiffeomorphism{}$\widehat{\varphi}\colon\widehat{M}\to\widehat{N}$ 给出，其中被微的那个函数%
\trCompactlySupported。特别地，这样的映射在一个紧子集之外保持\trLiouville\trVectorField{}$\partial_r$
不变。于是，对某个 $C\gg 0$，在%
$[C;\infty)\times\partial M\subset\widehat{M}$ 上它必形如
\begin{equation*}
\widehat{\varphi}(r,y)=\bigl(r-w(y),\psi(y)\bigr)
\end{equation*}
其中 $\psi\colon\partial M\to\partial N$ 是一个\trContact\trDiffeomorphism，而
$w\in C^\infty(\partial M,(-\infty;C])$ 是满足 $\psi^*(\theta_N|\partial N)=e^w(\theta_M|\partial M)$
的函数。我们说 $\widehat{\varphi}$ 是\textbf{在无穷附近\trContact{}型}（of contact type near infinity）
的。全体这类\trIsomorphism{}所成的空间 $\mathrm{Iso}(\widehat{M},\widehat{N})$ 带有一个自然的%
\trTopology；这容许那些在无穷远处改变\trContact\trDiffeomorphism、从而其\trSupport{}不\trCompact{}的%
\trIsotopy。照例，我们记 $\mathrm{Aut}(\widehat{M})=\mathrm{Iso}(\widehat{M},\widehat{M})$。一个标准
事实是：\trMoser{}引理的一个版本在此语境中成立：
\begin{lemma}\label{thm:contact-type-moser}
设 $(\varpi_M^s,\theta_M^s)$，$s\in[0;1]$，是 $M$ 上带\trContactTypeBoundary{}的\trExactSymplectic{}%
结构的一个\trSmooth{}族。那么，在所得的非紧\trManifold{}$\widehat{M}^0$ 与 $\widehat{M}^1$ 之间
存在一个在无穷附近\trContact{}型的\trExactSymplectic\trIsomorphism。
\noqed
\end{lemma}
设 $M$ 是\trContactTypeBoundary{}的\trExactSymplectic\trManifold。循着一个众所周知的策略（例如见
[76]），我们考虑 $\varpi_M$-\trCompatible{}的\trAlmostComplex{}结构 $J$，它在 $\partial M$ 的某个%
邻域内沿 $X_{\theta_M}$ 的\trFlow{}不变，并且满足
\begin{equation}\label{eq:contact-type-near-boundary}
dh_M\circ J=-\theta_M,
\end{equation}
其中 $h_M$ 如\eqref{eq:contact-gluing-function}。我们把这样的\trAlmostComplex{}结构称为%
\textbf{在边界附近\trContact{}型}（of contact type near the boundary）的。给定带%
\trComplexStructure{}$I_S$ 的\trRiemann\trSurface{}$S$，以及一个 $J$-\trHolomorphic{}映射
$u\colon S\to M$，考虑函数 $\rho=h_M\circ u\colon u^{-1}(U)\to(0;1]$。由
\eqref{eq:contact-type-near-boundary}，它满足
\begin{equation}\label{eq:rho-laplacian}
\Delta\rho=-d(d\rho\circ I_S)=-d(dh_M\circ J\circ du)=d(\theta_M\circ du)=u^*\varpi_M\ge 0.
\end{equation}
由\trMaximumPrinciple，若 $\rho$ 在 $S$ 的某个（内部）点处取到其最大可能值 $1$，则它在该处必然%
\trLocallyConstant。因此，任取一个在边界附近\trContact{}型的 $I_M=J$，便使 $\partial M$
弱\trConvex，从而把 $M$ 变成\trCorneredBoundary{}的\trExactSymplectic\trManifold。
我们将需要这个基本\trConvex{}性论证的若干变体。首先，它适用于\trAlmostComplex{}结构族，其意义如下：
\begin{lemma}\label{thm:family-convexity}
设 $(J^z)_{z\in S}$ 是 $M$ 上一族 $\varpi_M$-\trCompatible{}\trAlmostComplex{}结构，由某个%
\trRiemann\trSurface{}$S$ 所\trParametrization。假设每个 $J^z$ 都在 $\partial M$ 附近%
\trContact{}型。那么 $S\times M$ 上由 $J(z,x)=I_S(z)\times J^z(x)$ 给出的\trAlmostComplex{}结构
使 $S\times\partial M$ 弱\trConvex。
\noqed
\end{lemma}
其次，我们可以考察\eqref{eq:cone-completion} 型\trCompletion{}中的\trPseudoholomorphic{}曲线。$M$ 上每个
在边界附近\trContact{}型的\trAlmostComplex{}结构都有一个延拓到 $\widehat{M}$ 的\trCanonical{}%
$\hat{J}$，它在锥上对 $r$-\trTranslation{}不变。这满足
$dh_{\widehat{M}}\circ\hat{J}=-\theta_{\widehat{M}}$；施以与上文相同的计算\footnote{译注：原书此处
印作 “so by using the same computation as before. One gets:”，句号把一句话截断；疑应为
“…as before, one gets:”（同一句）；此处径改为一句直陈。——译者注}，即得：
\begin{lemma}\label{thm:completion-maximum-principle}
设 $S$ 是带边的\trCompact\trConnected\trRiemann\trSurface，$u\colon S\to\widehat{M}$ 是一个%
$\hat{J}$-\trHolomorphic{}映射，使得 $u(\partial S)\subset M$。那么 $u(S)\subset M$ 同样成立。
\noqed
\end{lemma}
还有一个积分版的\trMaximumPrinciple{}论证，它可用来给出上一引理的另一个证明，但也能产生稍更一般的%
结果。就是说，设 $N$ 是一个任意\trExactSymplectic\trManifold，$M$ 是\trContactTypeBoundary{}的%
\trExactSymplectic\trManifold（维数相同），而 $\iota\colon M\hookrightarrow\mathrm{int}(N)$ 是一个%
满足 $\varpi_M=\iota^*\varpi_N$、$\theta_M=\iota^*\theta_N$ 的\trEmbedding。为简便起见，我们把它
写作一个包含关系。取一个 $\varpi_N$-\trCompatible{}\trAlmostComplex{}结构 $J$，其在 $M$ 上的限制
在边界附近\trContact{}型。设 $S$ 是带边的\trCompact\trConnected\trRiemann\trSurface，
$u\colon S\to N\setminus\mathrm{int}(M)$ 是一个 $J$-\trHolomorphic{}映射，使得
$u(\partial S)\subset\partial M$。显然，
\begin{equation}\label{eq:integrated-maximum-principle}
0\le\int_S u^*\varpi_N=\int_{\partial S}u^*\theta_M=\int_{\partial S}u^*(-dh_M\circ J)
=\int_{\partial S}-dh_M\circ du\circ I_S.
\end{equation}
若 $\xi$ 是沿 $\partial S$ 的一个切向量，关于边界\trOrientation{}为正，则 $I_S\xi$ \trIncoming，
而 $du(I_S\xi)$ 不指向 $M$ 的内部；故 $dh_M(du(I_S\xi))\ge 0$。除非 $u$ 是常值，这与
\eqref{eq:integrated-maximum-principle} 矛盾。作为一个简单的应用，我们得到：
\begin{lemma}\label{thm:relative-maximum-principle}
设 $S$ 是带边的\trCompact\trConnected\trRiemann\trSurface，$u\colon S\to N$ 是一个
$J$-\trHolomorphic{}映射，使得 $u(\partial S)\subset\mathrm{int}(M)$。那么
$u(S)\subset\mathrm{int}(M)$ 同样成立。
\end{lemma}
\begin{proof}
把 $M$ 换成稍小一些的子集 $M^\epsilon=M\setminus h_M^{-1}((e^{-\epsilon};1])$，其中某个
$\epsilon>0$。通过适当地选取 $\epsilon$，我们可以做到：$J|M^\epsilon$ 仍在边界附近\trContact{}型、
$u(\partial S)$ 落在 $\mathrm{int}(M^\epsilon)$ 内，并且 $u$ 与 $\partial M^\epsilon$
\trTransverse{}相交。于是 $S^\epsilon=u^{-1}(N\setminus\mathrm{int}(M^\epsilon))$ 仍是一个带边的%
\trCompact\trRiemann\trSurface，其边界被映入 $\partial M^\epsilon$。把上一论证应用于每个%
\trConnectedComponent，我们发现 $u(S^\epsilon)\subset M^\epsilon$，从而
$u(S)\subset\mathrm{int}(M)$。
\end{proof}
最后，给定任意\trCorneredBoundary{}的\trExactSymplectic\trManifold{}$N$，可以按如下方式%
\emph{\trRoundOffCorners}。设 $F_1,\ldots,F_d$ 是 $N$ 的\trCodimension{}为 1 的\trBoundaryFace。
对每个 $F_j$ 选取\trSmooth{}函数 $f_j\colon N\to\RR^+$，使得 $F_j=f_j^{-1}(0)$，且 $df_j$ 在
$F_j$ 的外指法向上小于零，于是特别地 $X_{\theta_N}.f_j>0$ 在 $F_j$ 的某个邻域内成立。令
$f=f_1f_2\cdots f_d$。我们断言：在 $\mathrm{int}(N)$ 中离边界足够近的一切点处都有
$X_{\theta_N}.f>0$。这个不等式可以改写为
\begin{equation}\label{eq:rounding-inequality}
\sum_j\frac{X_{\theta_N}.f_j}{f_j}>0.
\end{equation}
设 $\mathrm{int}(N)$ 中有一个点列\trConvergence{}到一个极限点，使得该极限属于 $F_j$ 当且仅当
$j\in E$，这里 $E\subset\{1,\ldots,d\}$ 是某个非空子集。沿这个点列，
\eqref{eq:rounding-inequality} 中 $j\in E$ 的那些项各自趋于 $+\infty$，而其余的项\trBounded，
故该不等式终将成立。断言既已建立，立刻可见
\begin{lemma}\label{thm:rounding-lemma}
$M=f^{-1}([\delta;\infty))\subset N$，对任意充分小的 $\delta$（并赋以 $\varpi_N$、$\theta_N$
的限制），是\trContactTypeBoundary{}的\trExactSymplectic\trManifold。此外，
负\trLiouville\trFlow{}（\emph{取长时间}）会把整个 $N$ 压入 $\mathrm{int}(M)$ 之内。
\noqed
\end{lemma}
